\chapter{Complexity — Compile}
%%% generated by the blueprint pipeline from TCSlib.Complexity.SpaceComplexity.Machines.Compile
\section{Overview}

Declarations of the Lean module \texttt{TCSlib.Complexity.SpaceComplexity.Machines.Compile}.

\section{Declarations}

\begin{definition}[Finite tape word]
\lean{Complexity.LogProg.tapeWord}
\label{Complexity.LogProg.tapeWord}
\leanok
\uses{Turing.FinTM.bufferTape}
For a tape $f\colon\mathbb{Z}\to\mathrm{Option}(\{\mathrm{false},\mathrm{true}\})$,
define its tape word to be the unique finite Boolean word $w$ such that $f$ is the tape
encoding of $w$, if such a word exists; if no such word exists, define the tape word to
be the empty word.
\end{definition}

\begin{lemma}[Recovering a word from its buffer tape]
\lean{Complexity.LogProg.tapeWord_bufferTape}
\label{Complexity.LogProg.tapeWord_bufferTape}
\leanok
\uses{Complexity.LogProg.tapeWord, Turing.FinTM.bufferTape}
\difficulty{4}
For every finite Boolean word $w$, the word extracted from the tape storing $w$
contiguously from cell zero is $w$.
\end{lemma}

\begin{definition}[Register words in a configuration]
\lean{Complexity.LogProg.regWords}
\label{Complexity.LogProg.regWords}
\leanok
\uses{Complexity.LogProg.tapeWord, Turing.Cfg}
Given a configuration of an $m$-work-tape machine with Boolean work-tape symbols, and a
register index $r$, define the word in register $r$ to be the tape word of the contents
of work tape $r$.
\end{definition}

\begin{definition}[One step of a register program]
\lean{Complexity.LogProg.rstep}
\label{Complexity.LogProg.rstep}
\leanok
\uses{Complexity.LogProg.RProg, Complexity.LogProg.callSegs, Complexity.LogProg.regWords, Complexity.LogProg.vword, Turing.Cfg, Turing.MultiTapeTM.step}
Let $P$ be a register program, let an oracle assign a Boolean answer to each decider and
finite Boolean word, and let $c$ be a configuration on input $x$. Define one step as
follows: if $c$ has no state, leave it unchanged. If its state has no associated call,
take one ordinary transition of the program's underlying Turing machine; otherwise, form
the virtual word specified by that call, the input, and the selected register words,
query the oracle for the call's decider on that word, set the state to the call's
positive-answer state or negative-answer state according to the answer, and reset the
input-head position to $1$, leaving the other configuration fields unchanged.
\end{definition}

\begin{definition}[Repeated execution of a register program]
\lean{Complexity.LogProg.rrun}
\label{Complexity.LogProg.rrun}
\leanok
\uses{Complexity.LogProg.RProg, Complexity.LogProg.rstep, Turing.Cfg}
Let $P$ be a register program, let an oracle assign a Boolean answer to each decider and
finite Boolean word, let $c$ be a configuration on input $x$, and let $n\in\mathbb{N}$.
Define the configuration after $n$ steps to be the result of applying the program's
one-step transition $n$ times to $c$.
\end{definition}

\begin{lemma}[Successor recurrence for a run]
\lean{Complexity.LogProg.rrun_succ}
\label{Complexity.LogProg.rrun_succ}
\leanok
\uses{Complexity.LogProg.RProg, Complexity.LogProg.rrun, Complexity.LogProg.rstep, Turing.Cfg}
\difficulty{1}
Let $P$ be a register program with oracle, let $c$ be a configuration, and let
$n\in\mathbb N$. The run of length $n+1$ from $c$ is obtained by taking one step from
the configuration reached after $n$ steps.
\end{lemma}

\begin{definition}[Validity of oracle-call configurations]
\lean{Complexity.LogProg.CallOK}
\label{Complexity.LogProg.CallOK}
\leanok
\uses{Complexity.LogProg.CleanRun, Complexity.LogProg.RProg, Complexity.LogProg.callSegs, Complexity.LogProg.regWords, Complexity.LogProg.vword, Turing.Cfg, Turing.FinTM.bufferTape, Turing.MultiTapeTM}
Let $P$ be a register program, $D$ a decider with $k_D$ work tapes and state type $S_D$,
$q_0$ a choice of initial decider state for each of the $d$ deciders, and $O$ an oracle.
Let $\ell,r\colon\mathrm{Fin}(m)\to\mathbb{Z}$ be register bounds, let $B\in\mathbb{N}$,
and let $c$ be a configuration on input $x$. The configuration satisfies the call
conditions if, whenever its state is $\lambda$ and the program specifies a call at
$\lambda$, the call's argument list has no repetitions, the input head is at position
$1$, and each argument register contains exactly the tape encoding of its register word,
has its head at position $0$, and has word length between the bounds $\ell$ and $r$
(that is, $\ell$ is at most $-1$ and the word length is at most $r$). In addition, the
decider $D$, started in state $q_0$ for the specified decider on the call's virtual
word, satisfies the clean-run condition with bound $B$ and answer given by $O$ on that
word.
\end{definition}

\begin{theorem}[Correctness of the compiled simulation]
\lean{Complexity.LogProg.compile_correct}
\label{Complexity.LogProg.compile_correct}
\leanok
\uses{Complexity.LogProg.CallOK, Complexity.LogProg.RProg, Complexity.LogProg.compileTM, Complexity.LogProg.compile_correct_from, Complexity.LogProg.rrun, Complexity.LogProg.seam, Complexity.LogProg.seam_init, Turing.Cfg.init, Turing.MultiTapeTM, Turing.MultiTapeTM.initCfg, Turing.MultiTapeTM.runFrom}
\difficulty{1}
Let $P$ be a register program with $m$ registers and $d$ oracle deciders, let $D$ be a
decider machine with $k_D$ work tapes, and choose an initial program state, an initial
decider state for each oracle, an oracle, integer lower and upper register-head bounds
$lo,hi$, and a natural-number space bound $B$. Let $x$ be the input and $N$ a natural
number. Suppose that throughout the first $N+1$ configurations of the register-program
run from its initial configuration on $x$, every register head lies between its bounds.
Suppose also that at every time $t<N$, the current configuration satisfies the call
conditions: each call's argument list has no repetitions, the input head is at position
$1$, every argument register contains its represented word from cell $0$ with head at
$0$ and word length within the prescribed bounds, and the corresponding decider run on
the encoded virtual input halts with the oracle's answer, clears its work tapes, returns
their heads to $0$, and keeps every decider head within distance $B$ of $0$. Then the
compiled machine, started on $x$, reaches after some time $T$ exactly the compiled
configuration corresponding to the register-program configuration at time $N$. At every
time up to $T$, its original-register heads satisfy the bounds $lo,hi$, and each
decider-tape head has absolute position at most $B$.
\end{theorem}

\begin{theorem}[Space bound from head-position intervals]
\lean{Complexity.LogProg.spaceUsed_le_of_ranges}
\label{Complexity.LogProg.spaceUsed_le_of_ranges}
\leanok
\uses{Turing.Cfg, Turing.MultiTapeTM, Turing.MultiTapeTM.runFrom, Turing.MultiTapeTM.spaceUsed, Turing.MultiTapeTM.spaceUsedByTape, Turing.MultiTapeTM.visitedByTapeHead}
\difficulty{4}
Let $M$ be a deterministic machine with $K$ work tapes, let $c$ be an initial
configuration, and let $T\in\mathbb N$. Suppose that at every time $t\le T$, the head on
each tape $i$ lies between integers $\mathrm{lo}(i)$ and $\mathrm{hi}(i)$, inclusive.
Then the total number of distinct work-tape cells visited through time $T$ is at most
$\sum_i (\mathrm{hi}(i)-\mathrm{lo}(i)+1)$.
\end{theorem}

\begin{theorem}[Space bound for compiled computation]
\lean{Complexity.LogProg.compile_space}
\label{Complexity.LogProg.compile_space}
\leanok
\uses{Complexity.LogProg.CallOK, Complexity.LogProg.RProg, Complexity.LogProg.compileFinTM, Complexity.LogProg.compileTM, Complexity.LogProg.compile_correct, Complexity.LogProg.rrun, Complexity.LogProg.spaceUsed_le_of_ranges, Turing.Cfg.init, Turing.FinTM.ComputesInTime, Turing.FinTM.computesInTime_iff, Turing.MultiTapeTM, Turing.MultiTapeTM.initCfg, Turing.MultiTapeTM.runFrom, Turing.MultiTapeTM.spaceUsed}
\difficulty{5}
Let the program state set and decider state set be finite. Let $P$ be a register program
with $m$ registers and $d$ oracle deciders, and let $D$ be a decider machine with $k_D$
work tapes. Choose a start state, initial decider states, an oracle, register-head
bounds $lo,hi$, and natural numbers $B,N$; let $w$ be a bit word. Suppose the
register-program run on $x$ is halted at time $N$ with output $w$, its register heads
stay between $lo$ and $hi$ through time $N$, and at every earlier time its calls satisfy
the conditions that argument registers are distinct, the input head is at position $1$,
each argument register stores its represented word from cell $0$ with head at $0$ and
word length within the bounds, and the corresponding decider run halts with the oracle's
answer, clears its work tapes, returns their heads to $0$, and keeps each decider head
within distance $B$ of $0$. Then the compiled finite machine halts on $x$ with output
$w$ after some time $T$, and its total space use is at most $\sum_r (hi(r)-lo(r)+1) +
k_D(2B+1)$.
\end{theorem}

%%% added by the blueprint pipeline
\section{Additional declarations}

\begin{theorem}[Compiled simulation from an arbitrary configuration]
\lean{Complexity.LogProg.compile_correct_from}
\label{Complexity.LogProg.compile_correct_from}
\leanok
\uses{Complexity.LogProg.CallOK, Complexity.LogProg.RProg, Complexity.LogProg.callSegs, Complexity.LogProg.call_run, Complexity.LogProg.compileTM, Complexity.LogProg.regWords, Complexity.LogProg.rrun, Complexity.LogProg.rrun_succ, Complexity.LogProg.rstep, Complexity.LogProg.runFrom_forall_append, Complexity.LogProg.seam, Complexity.LogProg.step_seam_prog, Complexity.LogProg.vword, Turing.Cfg, Turing.MultiTapeTM, Turing.MultiTapeTM.runFrom, Turing.MultiTapeTM.runFrom_add, Turing.MultiTapeTM.runFrom_zero, Turing.MultiTapeTM.step}
\difficulty{6}
Let $P$ be a register program, let $D$ be a decider machine, and choose an initial
program state, initial decider states, an oracle, register-head bounds $lo,hi$, an
arbitrary initial configuration $c_0$, and a natural number $N$. Suppose every register
head in the run from $c_0$ stays between its bounds through time $N$, and every
configuration before time $N$ satisfies the call conditions: argument registers are
distinct, the input head is at position $1$, each argument register stores its
represented word from cell $0$ with head at $0$ and word length within the bounds, and
each corresponding decider run halts with the oracle's answer, clears its work tapes,
returns their heads to $0$, and keeps those heads within distance $B$ of $0$. Then the
compiled machine started in the configuration obtained by adjoining blank decider tapes
and compiled program state to $c_0$ reaches after some time $T$ the corresponding
compiled configuration for the register-program run at time $N$. At every time up to
$T$, its original-register heads remain within $lo,hi$, and each decider-tape head has
absolute position at most $B$.
\end{theorem}

\begin{definition}[Compilation into a finite Turing machine]
\lean{Complexity.LogProg.compileFinTM}
\label{Complexity.LogProg.compileFinTM}
\leanok
\uses{Complexity.LogProg.CSt, Complexity.LogProg.RProg, Complexity.LogProg.compileTM, Turing.FinTM, Turing.MultiTapeTM}
Let $P$ be a register program with $m$ registers and $d$ deciders, and let $D$ be a
decider with $k_D$ work tapes. Given an initial program state and a choice of initial
state for each decider, define the compiled finite Boolean Turing machine by compiling
these data into a machine with $m+k_D$ work tapes.
\end{definition}
